\chapter{Interfaces et captures d'écran}
\label{chap:captures}

Ce chapitre présente l'application Angular LogiFlow vue par un administrateur authentifié.
Les captures suivent le parcours d'usage : connexion et shell, tableau de bord, cycle
commande--dossier--voyage, flotte et exploitation, puis les deux usages principaux de l'IA.

\section{Connexion et App Shell}

L'entrée dans l'application passe par Keycloak (OIDC). Une fois authentifié, l'App Shell
fixe la navigation (sidebar filtrée par rôle) et les raccourcis globaux (Copilote, File du jour).

\figshot{01-connexion}{Écran de connexion : gateway vers Keycloak.}

\figshot{02-app-shell}{App Shell après connexion : modules, Copilote et File du jour visibles.}

\section{Tableau de bord}

Le tableau de bord synthétise les modules accessibles (compteurs), la file du jour
(exceptions prioritaires) et une répartition synthétique (ex.\ véhicules).

\figshot{03-tableau-de-bord}{Tableau de bord chargé : aperçu KPI, file du jour et synthèse.}

\section{Cycle commercial et transport}

Le cycle métier commence par les commandes, se poursuit par le dossier transportable,
puis le voyage (création assistée ou liste d'exécution avec carte OSRM).

\figshot{04-liste-commandes}{Liste des commandes (recherche et filtres).}

\figshot{05-fiche-dossier}{Fiche dossier de transport : segments, statut et marchandises.}

\figshot{06-nouveau-voyage}{Nouveau voyage : choix du mode planification assistée (IA).}

\figshot{07-liste-voyages}{Liste des voyages avec panneau d'itinéraire (OSRM).}

\section{Flotte, carburant et maintenance}

Les modules satellites couvrent le parc roulant, les prises de carburant et l'atelier.

\figshot{08-fiche-vehicule}{Fiche véhicule (ex.\ ME-001-CE) : identification, marque, poids et statut Disponible.}

\figshot{09-carburant}{Module carburant : prises et stations.}

\figshot{10-maintenance}{Module maintenance : ordres de travail et plans d'entretien.}

\section{Intelligence artificielle}

Deux écrans illustrent l'IA côté utilisateur : le copilote conversationnel (façade Spring)
et la planification assistée qui propose des voyages comparables.

\figshot{11-copilote}{Copilote en ligne : question métier et réponse structurée (voyages en cours).}

\figshot{12-planification-ia}{Planification assistée : propositions comparées et bouton de choix.}
